\clearpage
\fancyhead[L]{\small MERIDIANKREIS.}
\hypertarget{native-434-part-lower}{}
\begin{center}
{\Large MERIDIANKREIS.}

\rule{0.6\linewidth}{0.4pt}

Göttingische gelehrte Anzeigen. Stück 91. Seite 905..912. 1820. Juni 5.

\rule{0.6\linewidth}{0.4pt}
\end{center}

Der neue \textsc{Reichenbach}sche Meridiankreis, welchen unsere Sternwarte im vorigen Jahre erhalten
hat, ist zwar schon seit geraumer Zeit aufgestellt, und im Gebrauch: wir haben aber, ehe wir von
demselben eine öffentliche Nachricht gäben, erst eine beträchtliche Reihe von Erfahrungen abwarten
zu müssen geglaubt, um nach diesen desto sicherer über das, was dieses Instrument, von einer bisher
noch nicht angewandten Einrichtung, leistet, urtheilen zu können.

Das Instrument ist zugleich Mittagsfernrohr und Höhenmesser. Alle zu der erstern Bestimmung
gehörigen Einrichtungen hat es mit den vollkommensten einfachen Mittagsfernröhren gemein, und
bedürfen daher jene keiner ausführlichen Beschreibung. Das Fernrohr hat 5 Pariser Fuss Brennweite,
und 4 Pariser Zoll Oeffnung. Die vier Oculareinsätze vergrössern 68, 86, 120 und 170 mal: Hr.
Hofrath \textsc{Gauss} bedient sich fast ausschliesslich der stärksten. Das Fadennetz musste, weil
einige der Fäden schlaff geworden waren, hier erneuert werden. Es besteht jetzt aus sieben
verticalen und zwei horizontalen Spinnenfäden: die Zwischenräume zwischen jenen werden von Sternen
im Aequator in $\cdot 14$ Secunden durchlaufen; die horizontalen sind nur $7''6$ von einander
entfernt.

Die Axe von 33 Pariser Zoll Länge trägt auf der einen Seite zwei concentrische Kreise, deren vom
Fernrohr abgekehrte Flächen fast in Einer Ebene liegen. Der äussere Kreis, welcher auf der Axe
\textit{fest} ist, und sich also mit dem Fernrohr zugleich dreht, hat die Theilung, die unmittelbar
von 3 zu 3 \SourceLeaf{435}
Minuten geht. Der innere Kreis (Alhidaden-Kreis) würde sich ohne eine an dem Pfeiler angebrachte
Hemmung um die Axe drehen lassen: diese aber verstattet ihm nur eine kleine feine Bewegung, um die
an ihm feste Wasserwaage einzustellen. Auf diesem Alhidadenkreise sind die vier Indices, je
$45^\circ$ von der Verticallinie entfernt, mit ihren Verniers, welche die Haupttheile wieder in 90
Theile, also von 2 zu 2 Secunden abtheilen; kleinere Theile lassen sich noch schätzen. Der
Durchmesser der Kreise bei der Ablesung ist 35 Pariser Zoll. Dass beide Kreise, ohne einander zu
berühren, doch nur durch einen kaum merklichen Zwischenraum getrennt, und dass desswegen die
Ablesungsmicroscope absichtlich etwas schief gestellt sind, indem die Fläche des getheilten Kreises
etwas, obwohl nur äusserst wenig, über die Fläche des Alhidadenkreises hervorragt, sind
Einrichtungen, die dieses Instrument mit andern \textsc{Reichenbach}schen gemein hat. Drei
Balancirungen, für das ganze Instrument, für den Alhidadenkreis und für das Fernrohr, haben den
Zweck, den Druck der Zapfen auf die Pfannen und des Alhidadenkreises auf die Axe, so wie die
Biegung des Rohrs aufzuheben. An der Hänglibelle, womit die Axe nivellirt wird, beträgt ein
Ausschlag von einem Pariser Zoll $22''$, an der Hauptlibelle $17''6$. Letztere dient, den
Alhidadenkreis immer in unverrückter Lage zu erhalten, oder sehr kleine Verrückungen zu messen und
in Rechnung zu bringen (Hr. Hofr. \textsc{Gauss} pflegt die Correctionsschraube nur dann zu
berühren, wenn die Verstellung über $2''$ beträgt, welches selten der Fall ist).

Die Umlegung des Instruments, die vermittelst einer sehr zweckmässig eingerichteten Maschine
leicht und sicher ausgeführt wird, dient demnächst, in Beziehung auf das Höhenmessen, den
Collimationsfehler auszumitteln, und die gemessenen Zenithdistanzen in absolute zu verwandeln.
Inzwischen kann es bedenklich scheinen, den Collimationsfehler auf längere Zeit und bei
beträchtlichem Temperaturwechsel als ganz unveränderlich anzusehen, und daher wird es, falls man
nicht das Instrument sehr oft umlegt, vortheilhafter sein, die Beobachtungen auf den Pol, als sie
auf das Zenith zu beziehen. Den Ort des Pols auf dem Kreise kann man durch Circumpolarsterne um so
leichter bestimmen, da die optische Kraft des Fernrohrs bewundernswürdig gross ist. Dieses steht
darin dem Fernrohr am \textsc{Reichenbach}schen Passage-Instrument (d. Anz. 1819 [S. 422 d. B.])
kaum nach, obgleich letzteres bei einer grössern Brennweite auch eine etwas grössere Oeffnung hat,
und der etwa noch vorhandene kleine Unterschied wird dadurch aufgewogen, dass man den
Meridiankreis vorher auf das genaueste einstellen kann, und so allemal seine ganze Aufmerksamkeit
nur auf einen sehr kleinen Theil des Gesichtsfeldes zu richten braucht. Unter den von Hrn. Hofrath
\textsc{Gauss} bei Tage beobachteten Sternen finden sich eine Menge von der vierten Grösse; 1 Draconis \textsc{Hevelii} und $\omega$ Cephei \textsc{Hevelii}, welche auch bei Tage und in der
Nähe des Mittags beobachtet wurden, sind sogar nur von der fünften. Die Einrichtung des
Fadennetzes bringt es aber mit sich, dass die Tagesbeobachtungen so feiner Pünktchen, sobald sie
nur überhaupt sichtbar sind, dieselbe Zuverlässigkeit haben, wie die der grössern Sterne.

\Needspace{40\baselineskip}
Wir setzen hier als eine Probe der herrlichen mit diesem Instrument zu erreichenden
Uebereinstimmung die Resultate aller von Eröffnung des Tagebuches bis zur ersten Umlegung gemachten
Bestimmungen des Orts des Pols auf dem Instrumente, aus allen unmittelbar nach einander beobachteten
entgegengesetzten Culminationen von Circumpolarsternen her. Die Harmonie der Resultate aus einem und
demselben Sterne beweist die Feinheit, womit sich die Beobachtungen anstellen lassen, und die
während jener Zeit bewahrte Unveränderlichkeit des Collimationsfehlers: die eben so grosse
Uebereinstimmung der Resultate aus verschiedenen Sternen hingegen beweist die Vortrefflichkeit der
\textsc{Reichenbach}schen Theilung, die um so mehr Bewunderung verdient, wenn man die mässige
Grösse des Kreises in Betrachtung zieht.

\Needspace{29\baselineskip}
\SourceLeaf{436}
\begin{center}
\begin{tabular}{l|l|c|c}
1820 & Name der Sterne & Grösse & Ort des Pols \\
\hline
Febr. 26 & $\alpha$ Cephei & 3 & $321^\circ\,29'\,32''14$ \\
 & $\beta$ Cephei & 3 & $31.89$ \\
 & $\delta$ Draconis & 3 & $32.29$ \\
 & $\eta$ Cephei & 3.4 & $32.24$ \\
 & $\alpha$ Cephei & 3 & $31.16$ \\
 & $\beta$ Cephei & 3 & $31.65$ \\
27 & $\delta$ Draconis & 3 & $31.39$ \\
 & $\eta$ Cephei & 3.4 & $31.90$ \\
 & $\delta$ Draconis & 3 & $31.54$ \\
 & $\eta$ Cephei & 3.4 & $31.63$ \\
28 & $\gamma$ Cephei & 3 & $31.54$ \\
 & $\tau$ Cephei & 4 & $32.74$ \\
 & $\delta$ Cephei & 4 & $31.12$ \\
März 9 & $\delta$ Draconis & 3 & $31.18$ \\
 & $\beta$ Cephei & 3 & $32.18$ \\
 & $\tau$ Cephei & 4 & $31.42$ \\
 & $\iota$ Cephei & 4 & $31.52$ \\
10 & $\beta$ Cephei & 3 & $32.05$ \\
11 & 1 Dracon. Hev. & 5 & $31.97$ \\
18 & 1 Dracon. Hev. & 5 & $32.04$ \\
 & 3 Lacertae & 4 & $32.32$ \\
 & $\alpha$ Ursae maj. & 2 & $31.62$ \\
19 & $\alpha$ Ursae min. & 2 & $32.81$ \\
\end{tabular}
\end{center}

\par\indent Das Mittel ist $321^\circ\,29'\,31''84$, von welchem auch nicht eine einzige dieser 23
Bestimmungen eine volle Secunde abweicht.

\par\indent In so fern die Beobachtungen an diesem Instrument nicht auf das Zenith, sondern auf den Pol bezogen
werden, ist im Grunde die genaueste Kenntniss der Polhöhe etwas untergeordnetes, und da man, wo es
die einzelne Secunde oder gar Theile von Secunden gilt, die grösste Behutsamkeit anzuwenden hat, und
eine \textit{so grosse} Unveränderlichkeit des Collimationsfehlers wie die obige Uebersicht zeigt,
doch nicht auf immer zu erwarten ist, so wird es für die schärfste Bestimmung der Polhöhe
zweckmässig sein, nur solche Beobachtungen zu verbinden, die dem Umlegen zunächst vorhergingen und
zunächst darauf folgten. Hr. Hofr. \textsc{Gauss} wagt daher nicht, seine bisherigen Beobachtungen
zu \textit{dieser} Bestimmung anzuwenden, da nach dem ersten Umlegen ziemlich lange anhaltendes
ungünstiges Wetter eintrat, und beim zweiten am 13. April ein zufälligerweise vorgekommenes kleines
Derangement des Fadennetzes die \textit{völlige} Gleichheit des Collimationsfehlers etwas
zweifelhaft machte, obwohl übrigens die Resultate aus jenen Beobachtungen sehr nahe unter sich und
mit einem auf einem andern Wege erhaltenen, welches wir sogleich anführen wollen, übereinstimmen.
Die Verbindung des Orts des Pols im ersten Zeitraum, mit dem im zweiten ($38^\circ\,25'\,53''\!54$)
gibt nämlich die Polhöhe $51^\circ\,31'\,49''\!15$; und die Verbindung von 13
Circumpolarsternbeobachtungen unmittelbar vor dem zweiten Umlegen mit den unmittelbar nach dem
Umlegen beobachteten entgegengesetzten Culminationen derselben Sterne gibt
$51^\circ\,31'\,47''\!86$.

Bekanntlich hat seit einigen Jahren ein Umstand in einem hohen Grade die Aufmerksamkeit der
Astronomen auf sich gezogen, der, wenn gleich es dabei nur ein paar Secunden gilt, doch für die
Beobachtungskunst sowohl als für die mannichfaltigen astronomischen Resultate von höchster
Wichtigkeit ist. Wir meinen die kleinen Unterschiede, welche sich in Bestimmung von
Sterndeclinationen, Schiefe der Ekliptik und Polhöhen bei Anwendung verschiedener, wenn gleich
übrigens höchst vortrefflicher, Instrumente gezeigt haben. Es leidet keinen Zweifel, dass diese
Unterschiede von der \emph{Einwirkung der Schwerkraft} auf die verschiedenen Theile jedes
Instruments herrühren, wenn man gleich bis \SourceLeaf{437}
jetzt weder die Art der Einwirkung vollständig und mit Gewissheit nachweisen, noch bestimmt
urtheilen kann, welches Instrument das richtige und welches das unrichtige Resultat gegeben habe.
Ueber das Quantitative der Biegung der Metalle wissen wir im Grunde noch sehr wenig, und es scheint
uns zu gewagt, bei irgend einem Instrumente, es sei gebaut wie es wolle, die Möglichkeit eines
merklichen Einflusses der Schwere auf die Theile desselben und dadurch auf die Beobachtungen
abzuläugnen, ohne eine solche Behauptung durch hinlängliche anderweitige Beweise zu begründen. Bei
dem hiesigen Meridiankreise hat zwar der grosse Künstler alles gethan, um die Biegung des Fernrohrs
durch eine sinnreiche Balancirung aufzuheben. Inzwischen könnte man doch immer noch zweifeln, ob
dadurch alle Flexion weggeschafft, oder in ihren Wirkungen vollkommen unmerklich gemacht sei. Das
einzige directe Mittel, dieses zu prüfen, scheint die Verbindung unmittelbarer Beobachtungen eines
Himmelskörpers mit solchen zu sein, die an dem aus einem künstlichen Horizont gesehenen Bilde
desselben gemacht werden, wobei sich von selbst versteht, dass dergleichen Beobachtungen oftmals
wiederholt werden müssen, um diesen delicaten Gegenstand ins Klare zu bringen. Einen Anfang mit
solchen Beobachtungen hat Hr. Hofr. \textsc{Gauss} bereits gemacht, indem er den Nordstern aus
einem Wasserspiegel beobachtete: den glänzendsten Beweis von der erstaunlichen optischen Kraft des
Fernrohrs gibt wohl der Umstand, dass auch die bei Tage einfallende obere Culmination auf diese
Weise recht gut beobachtet werden konnte. Das Resultat der ersten \textit{vollständigen}
Beobachtung dieser Art war folgendes:

\Needspace{22\baselineskip}
{\rightskip=0pt plus 3em \textit{Zenithdistanz des Nordsterns} (frei von Strahlenbrechung, aber einschliesslich des
Culminationsfehlers)\par}

\begin{center}
1820 Mai 13.
\end{center}

\begin{center}
\begin{tabular}{l c l}
Untere Culmination, direct & . . . . . . . . . . & $319^\circ\,50'\,20''\!73$ \\
Untere Culmination, reflectirt & . . . . . . . . . & $220\ \ 5\ \ 3.94$ \\
Obere Culmination, direct & . . . . . . . . . . & $323\ \ 8\ 41.51$ \\
Obere Culmination, reflectirt & . . . . . . . . . & $216\ 46\ 44.31$ \\
\end{tabular}
\end{center}

\begin{center}
Hieraus folgt die \textit{wahre} Zenithdistanz
\end{center}

\begin{center}
\begin{tabular}{l c l}
in der untern Culmination & . . . . . . . . . . & $40\ \ 7\ 21.60$ \\
in der obern Culmination & . . . . . . . . . . & $36\ 49\ \ 1.40$ \\
\end{tabular}
\end{center}

\noindent und hieraus, indem die 12stündige Aenderung der Declination $-0''10$ ist, die Polhöhe für den
Platz des Wassergefässes $51^\circ\,31'\,48''\!45$, oder für den Mittelpunkt des Kreises
$51^\circ\,31'\,48''\!40$.

\par\indent Da diese Bestimmung fast mitten zwischen die zwei oben angeführten fällt, so wird es hiedurch
bereits sehr wahrscheinlich, dass die Wirkung der Schwere auf die Beobachtungen mit diesem
Instrument entweder völlig unmerklich, oder doch äusserst gering sein muss. Die Declination des
Nordsterns aus obigen Beobachtungen ist $0''42$ kleiner als aus \textsc{Bessels} Tafeln (die
sämmtlichen bisherigen Beobachtungen des Hrn. Hofr. \textsc{Gauss} geben diese Correction $=
-0''67$).

\Needspace{17\baselineskip}
Um auch noch eine Probe von dem, was das Instrument als Mittagsfernrohr leistet, zu geben, setzen
wir hier noch einige Beobachtungen des Mars her, welche Hr. Hofr. \textsc{Gauss} um die Zeit der
Quadratur dieses Planeten gemacht hat.

\begin{center}
\begin{tabular}{l|c|c|c}
1820 & Mittlere Zeit & \shortstack{G. Aufst.\\d. 1. Rand.} & \shortstack{Abw. d.\\Mittelp.} \\
\hline
März 29 & $7^{\mathrm u}\,10^{\mathrm m}\,16^{\mathrm s}3$ & $114^\circ\,37'\,49''3$ & $24^\circ\,16'\,47''3$ N. \\
30 & $7\,7\,40.7$ & $114\,57\,57.7$ & $24\,12\,7.1$ \\
31 & $7\,5\,6.8$ & $115\,18\,30.1$ & $24\,7\,18.1$ \\
April 5 & $6\,52\,37.3$ & $117\,6\,19.9$ & $23\,42\,5.4$ \\
10 & $6\,40\,40.4$ & $119\,2\,17.5$ & $23\,14\,30.6$ \\
11 & $6\,38\,20\,3$ & $119\,26\,18.3$ & $23\,8\,41.0$ \\
12 & $6\,36\,1.6$ & $119\,50\,41.5$ & $23\,2\,48.0$ \\
13 & $6\,33\,43.7$ & $120\,15\,15.0$ & $22\,56\,44.8$ \\
26 & $6\,5\,16.8$ & $125\,56\,10.0$ & $21\,29\,12.8$ \\
\end{tabular}
\end{center}

\Needspace{18\baselineskip}
\SourceLeaf{438}
Die Vergleichung dieser Beobachtungen mit den Marstafeln des Hrn. von \textsc{Lindenau}
wurde durch Hrn. von \textsc{Staudt} gemacht, welcher sich gegenwärtig bei uns mit
ausgezeichnetem Eifer und Erfolg dem Studium der mathematischen und astronomischen
Wissenschaften widmet. Es ergaben sich folgende

\begin{center}
Unterschiede

\begin{tabular}{ll|r|r}
 & & Ger. Aufst. & Abweich. \\
\hline
März & 29 & $-0''7$ & $+5''1$ \\
 & 30 & $-0.5$ & $+3.9$ \\
 & 31 & $-2.1$ & $+6.1$ \\
April & 5 & $-1.3$ & $+3.9$ \\
 & 10 & $-4.9$ & $+2.9$ \\
 & 11 & $-1.9$ & $+3.6$ \\
 & 12 & $-4.8$ & $+1.9$ \\
 & 13 & $-1.8$ & $+4.2$ \\
 & 26 & $-1.5$ & $+1.5$ \\
\end{tabular}
\end{center}
